\begin{explanationbox}
	\textbf{\textcolor{CExplanation}{一句话直觉}}
	数列的单调性，可以通过相邻两项的差或商是否保持同一符号来判断。

	\medskip
	\textbf{\textcolor{CExplanation}{核心拆解}}
	\begin{description}[style=nextline, leftmargin=2.8em, labelsep=0.5em, itemsep=0.45em]
		\item[\textbf{要点一（差比法通用）：}]
			对任意数列，$a_{n+1}-a_n>0$等价于严格递增，$=0$等价于常数列，$<0$等价于严格递减。无需其他条件。
		\item[\textbf{要点二（商比法条件）：}]
			仅当所有$a_n>0$时，$a_{n+1}/a_n>1$等价于严格递增，否则不能直接使用。
	\end{description}

	\medskip
	\textbf{\textcolor{CExplanation}{几何本质（动点纵差）}}
	将数列$\{a_n\}$视为自然数$n$上的函数$f(n)=a_n$，则$a_{n+1}-a_n$就是离散的增量；其正负对应函数上升或下降。

	\medskip
	\textbf{\textcolor{CExplanation}{代数意义（差商转化）}}
	当$a_n>0$时，比较$a_{n+1}/a_n$与$1$等价于比较$a_{n+1}-a_n$与$0$（因为$a_n>0$不改变不等号方向），因此两种手段本质相同。

	\medskip
	\textbf{\textcolor{CExplanation}{考点价值（数列工具）}}
	\begin{itemize}[leftmargin=*, itemsep=0.5em, topsep=0.4em]
		\item\textbf{考法一（直接判定）：}给出通项或递推，要求判断数列的单调性。直接计算差或商并判断符号。\item\textbf{考法二（参数调控）：}已知数列含参数，且要求单调递增（减），利用差比法转化为恒成立问题，求出参数范围。
	\end{itemize}

	\medskip
	\textbf{\textcolor{CExplanation}{顿悟点（差商互换）}}
	差比法适合任意数列，商比法适合正项数列；遇到参数或含绝对值时优先使用差比法，正项且乘积形式时商比法往往更简便。

	\medskip
	\textbf{\textcolor{CExplanation}{使用场景}}
	\begin{itemize}[leftmargin=*, itemsep=0.5em, topsep=0.4em]
		\item\textbf{场景一（递推单调证明）：}已知正项递推式$a_{n+1}=f(a_n)$，直接计算$\frac{a_{n+1}}{a_n}$与$1$比较，快速得到单调性。\item\textbf{场景二（不等式放缩）：}利用数列单调性，可以确定项的范围，进而进行放缩或求和。
	\end{itemize}

\end{explanationbox}
